\subsection{Taxonomy normalization and mapping}
\label{methods:taxonomy}


A taxonomic label can denote different taxa across databases and
releases, and a taxon can have several labels \cite{balvovciute2017silva}.
We therefore retain the original SILVA lineage alongside a validated
NCBI Taxonomy identifier or a lineage-scoped project identifier.
Taxonomic assignments come from SILVA v138.2 (Section~\ref{methods}), and we treat the original classifier lineage and ASV identifier as the primary evidence.
These assignments are seven rank fields (domain--species), and some are followed by a separate species assignment field and a numeric confidence. We assert the species from the ranked lineage and retain the additional species call in a provenance ledger.
We normalize the seven ordered rank 
fields by removing rank prefixes and representing empty and unclassified entries explicitly.

The mapping procedure evaluates the historical SILVA-to-NCBI candidates against a pinned NCBI Taxonomy release (the OBO Foundry NCBITaxon release \texttt{2025-03-13}, built from the NCBI
organismal taxonomy of 2025-03-01).
An identifier is retained only when it resolves, has a compatible rank and label or documented synonym, and its ancestor chain is compatible with all retained
higher-rank identifiers in that source lineage. A missing, ambiguous
or incompatible candidate is replaced by a deterministic project class
scoped to the complete truncated source lineage and rank. Each decision records the source
lineage and rank, candidate identifier, validation outcome, final
identifier and reason. The output mapping covers every 
lineage represented in both the taxonomy export and the
feature table; otherwise generation stops.

Existing project identifiers are retained only when all explicit
historical uses share one normalized rank, label and nearest explicit
parent context; a reused identifier that acquires more than one
corrected context is replaced in every affected use by contextual
classes. Each resulting class retains one rank, label and parent context. The mapping builder emits
corrected mappings, a decision ledger, a clean project taxonomy module
and a manifest containing input and output checksums, and the
abundance-ABox generator accepts only a passing, checksum-matched
manifest.

Queries retain the original SILVA lineage and independently validated
NCBI or project identifiers. The decision ledger reports rejected
candidates and their reasons. Algorithmic checks establish consistency
with the specified name, rank and ancestry rules; semantic mapping
accuracy requires a separately annotated validation sample.
